\begin{quote}
\emph{Step scaling expansion with uniform remainder and $b_0>0$ in the chosen gradient-flow scheme.}
\end{quote}

\begin{enumerate}
\item It closes the \emph{quadratic-coefficient} part: the coefficient $b_0$ of the $u^2\log 2$ term in step scaling (equivalently the $u^2$ term in the beta function) is scheme invariant.
In the present submission we only use the fact that $b_0>0$ in the chosen GF scheme; a nonperturbative derivation of this positivity is given in the embedded Appendix~F (curvature localization + Brascamp--Lieb + cumulant expansion).

\item It reduces the \emph{uniform remainder} part to one explicit analytic lemma in our RG topology: a uniform (in $a\to 0$) bound
\begin{equation}
\label{eq:beta_remainder_goal}
\bigl|\beta_{\mathrm{GF}}(u) + 2 b_0 u^2\bigr|\le B\,u^3\qquad (0\le u\le u_0)
\end{equation}
for some $u_0,B>0$ independent of $a/L$. Once \eqref{eq:beta_remainder_goal} holds, the desired step-scaling remainder
$|R_\Sigma(u)|\le C\,u^3$ follows with completely explicit constants.
\end{enumerate}

\noindent\textbf{Note (load-bearing usage of $b_0$).}
In the load-bearing chain of this submission we do \emph{not} use the closed form of $b_0$.
We only use that the quadratic coefficient in the weak-window step-scaling/beta expansion is strictly positive in the chosen GF scheme.
That positivity is obtained by the repaired curvature-sector package in Appendix~\ref{app:route1};
see Theorem~\ref{thm:b0_coefficient_extraction} and Lemma~\ref{lem:b0_positive}, together with the downstream sufficiency certificate, Theorem~\ref{thm:weak_window_sufficiency}.
The classical closed form (the usual one-loop coefficient) is recorded only as an optional consistency check.

Fix a compact simple gauge group $G$ and a finite Euclidean box of linear size $L$ (either periodic $L^4$ or Schr\"odinger functional $L^3\times T$ with $\rho=T/L$ fixed). We work at bare lattice spacing $a$ with $L/a\in\mathbb N$.

Let $U$ be a lattice gauge field and $V_t(U)$ its Wilson flow at flow time $t\ge 0$.
Let $G_{\mu\nu}(t,x)$ be the field strength of the flowed field and define the flowed energy density
\begin{equation}
E(t,x):=\frac14\,\mathrm{tr}\,G_{\mu\nu}(t,x)G_{\mu\nu}(t,x).
\end{equation}

Fix parameters
\begin{equation}
 c\in(0,1),\qquad \rho\in(0,\infty),\qquad \xi:=x_0/T\in(0,1)
\end{equation}
(if Schr\"odinger functional) and set
\begin{equation}
\label{eq:t_def}
 t:=\frac{c^2L^2}{8},\qquad \mu:=\frac{1}{\sqrt{8t}}=\frac{1}{cL}.
\end{equation}
Define the renormalized coupling in the GF scheme by
\begin{equation}
\label{eq:gf_coupling_def}
 u_{\mathrm{GF}}(L)\equiv g^2_{\mathrm{GF}}(L):=\mathcal N^{-1}(c,\rho,\xi)\,t^2\,\langle E(t,x_0)\rangle_{L,T},
\end{equation}
where $\mathcal N(c,\rho,\xi)>0$ is the (tree-level) normalization chosen so that
\begin{equation}
\label{eq:tree_level_norm}
 g^2_{\mathrm{GF}}(L)=g_0^2+O(g_0^4)\qquad (g_0\to 0).
\end{equation}


\paragraph{Remark (normalization in infinite volume; optional).}
In infinite volume, a direct free-field (tree-level) computation fixes the leading normalization of
$t^2\langle E(t)\rangle$ and motivates the canonical choice of $\mathcal N$.
No explicit closed form is used below; only the tree-level matching property \eqref{eq:tree_level_norm} is used.

Define the (continuum) step scaling map at factor $2$ by
\begin{equation}
\label{eq:step_scaling_def}
\Sigma(u):=u_{\mathrm{GF}}(2L)\ \ \text{evaluated at bare coupling tuned so that }u_{\mathrm{GF}}(L)=u,
\end{equation}
followed by the continuum limit $a/L\to 0$ (holding $c,\rho,\xi$ fixed).

Let $u_{\mathrm{GF}}(\mu)$ denote the coupling at renormalization scale $\mu=1/(cL)$.
Define
\begin{equation}
\label{eq:beta_def}
\beta_{\mathrm{GF}}(u):=\mu\frac{du}{d\mu}=-\frac{d u}{d\ln L}.
\end{equation}
For asymptotically free Yang--Mills, $\beta_{\mathrm{GF}}(u)<0$ for small $u$.

For fixed factor $s>1$, step scaling is the time-$\ln s$ map of the ODE
\begin{equation}
\label{eq:rg_ode}
\frac{d u}{d\tau} = -\beta_{\mathrm{GF}}(u),\qquad \tau:=\ln L.
\end{equation}
In particular, $\Sigma(u)$ is obtained by flowing \eqref{eq:rg_ode} from $\tau$ to $\tau+\ln 2$.

\begin{proposition}[One-loop coefficient is scheme invariant]
\label{prop:b0_scheme_invariant}
Let $u$ and $u'$ be two renormalized couplings related near $0$ by an analytic reparametrization
\begin{equation}
\label{eq:scheme_change}
 u' = u + a_1 u^2 + a_2 u^3 + O(u^4)
\end{equation}
for some real coefficients $a_1,a_2$.
Let $\beta(u)=\mu\,du/d\mu$ and $\beta'(u')=\mu\,du'/d\mu$.
If
\begin{equation}
\label{eq:beta_expand}
 \beta(u) = -2 b_0 u^2 + O(u^3)
\end{equation}
then
\begin{equation}
\beta'(u') = -2 b_0 (u')^2 + O\bigl((u')^3\bigr).
\end{equation}
In particular, the coefficient $b_0$ is invariant under \eqref{eq:scheme_change}.
\end{proposition}

\begin{proof}
Differentiate \eqref{eq:scheme_change} with respect to $\ln\mu$:
$\beta'(u') = \frac{d u'}{d u}\,\beta(u)$.
From \eqref{eq:scheme_change},
$\frac{d u'}{d u}=1+2a_1u+3a_2u^2+O(u^3)$.
Insert \eqref{eq:beta_expand}:
\[
\beta'(u')=\bigl(1+2a_1u+O(u^2)\bigr)\bigl(-2 b_0 u^2+O(u^3)\bigr)=-2 b_0 u^2+O(u^3).
\]
Finally invert \eqref{eq:scheme_change} to write $u=u'+O((u')^2)$, so $u^2=(u')^2+O((u')^3)$, and the claim follows.
\end{proof}

\paragraph{Remark (orientation: perturbative matching; not used).}
In the conventional continuum treatment, one may match the gradient-flow coupling $u_{\mathrm{GF}}$ to a reference coupling
(such as $\overline{\mathrm{MS}}$) and compute the corresponding one-loop coefficient and higher-order terms.
In particular one obtains an expansion of the form
\[
\beta_{\mathrm{GF}}(u)=-2b_0u^2+O(u^3)\qquad (u\downarrow 0),
\]
with the usual one-loop coefficient $b_0=\frac{11}{3}C_2(\mathrm{ad})/(16\pi^2)$ in standard normalization.
\emph{None of this perturbative matching input is used in the load-bearing arguments below.}
Instead, in this submission:
(i) the strict positivity $b_0>0$ in the chosen GF scheme is obtained nonperturbatively along the $t_0$-defined scaling trajectory
(Appendix~F), and
(ii) the uniform remainder bound \eqref{eq:beta_remainder_goal} is obtained from the analytic dyadic RG estimates
(Theorem~\ref{thm:RG_analytic_bounds}) in the RG seminorm topology.

It remains to obtain a \emph{uniform} remainder bound
\begin{equation}
\label{eq:step_remainder_goal}
\Sigma(u)=u+2b_0u^2\ln 2 + R_\Sigma(u),\qquad |R_\Sigma(u)|\le C_\Sigma\,u^3,
\end{equation}
for $0\le u\le u_0$, with $u_0,C_\Sigma$ independent of $a/L$.

The next lemma shows that \eqref{eq:step_remainder_goal} follows from a single beta-remainder bound \eqref{eq:beta_remainder_goal}.

\begin{lemma}[From beta-remainder to step-scaling remainder with explicit constants]
\label{lem:beta_to_step_remainder}
Assume:
\begin{enumerate}
\item There exist constants $u_0>0$ and $B\ge 0$ such that the beta function satisfies
\begin{equation}
\label{eq:beta_remainder_assumption}
\bigl|\beta_{\mathrm{GF}}(u)+2b_0u^2\bigr|\le B u^3\qquad (0\le u\le u_0).
\end{equation}
\item $u_0$ is small enough that
\begin{equation}
\label{eq:u0_smallness}
D:=1-2b_0u_0\ln 2\ >\ 0.
\end{equation}
\end{enumerate}
Then for every $u\in[0,u_0]$, the factor-$2$ step scaling map satisfies
\begin{equation}
\label{eq:step_expansion_bound}
\Sigma(u)=u+2b_0u^2\ln 2 + R_\Sigma(u),
\qquad |R_\Sigma(u)|\le C_\Sigma\,u^3,
\end{equation}
with the explicit constant
\begin{equation}
\label{eq:Csigma}
C_\Sigma := \frac{(2b_0\ln 2)^2}{D}\; +\; \frac{B\ln 2}{D^3}.
\end{equation}
\end{lemma}

\begin{proof}
Let $u(\tau)$ solve the RG ODE \eqref{eq:rg_ode} with initial data $u(0)=u$.
Then $\Sigma(u)=u(\ln 2)$.
Write the ODE as
\begin{equation}
\label{eq:ode_split}
\frac{du}{d\tau} = 2b_0u^2 + r(u),\qquad r(u):= -\beta_{\mathrm{GF}}(u)-2b_0u^2.
\end{equation}
By \eqref{eq:beta_remainder_assumption}, $|r(u)|\le Bu^3$ for $u\le u_0$.

Define the comparison (exact one-loop) solution
\begin{equation}
\label{eq:w_def}
 w(\tau):=\frac{u}{1-2b_0u\tau},\qquad \text{so }\; w'(\tau)=2b_0 w(\tau)^2,
\end{equation}
which is well-defined on $[0,\ln 2]$ by \eqref{eq:u0_smallness}.

\paragraph{Step 1: bound the pure one-loop truncation error.}
A direct computation gives
\[
 w(\ln 2)=\frac{u}{1-2b_0u\ln 2}=u+2b_0u^2\ln 2+\frac{(2b_0\ln 2)^2}{1-2b_0u\ln 2}\,u^3.
\]
Therefore, for $u\le u_0$,
\begin{equation}
\label{eq:w_trunc_bound}
\bigl|w(\ln 2)-u-2b_0u^2\ln 2\bigr|\le \frac{(2b_0\ln 2)^2}{D}\,u^3.
\end{equation}

\paragraph{Step 2: control the perturbation by $r$.}
Let $\delta(\tau):=u(\tau)-w(\tau)$.
Subtract the ODEs \eqref{eq:ode_split} and \eqref{eq:w_def}:
\begin{equation}
\label{eq:delta_eq}
\delta'(\tau)=2b_0\bigl(u(\tau)^2-w(\tau)^2\bigr)+r\bigl(u(\tau)\bigr)=2b_0\delta(\tau)\bigl(u(\tau)+w(\tau)\bigr)+r\bigl(u(\tau)\bigr).
\end{equation}
Since $u(\tau)$ is increasing for small $u$ (asymptotic freedom), and $w(\tau)\le u/D$ on $[0,\ln 2]$, a bootstrap on \eqref{eq:delta_eq} yields
$u(\tau)\le u/D$ for $\tau\in[0,\ln 2]$ as well.
Hence $u(\tau)+w(\tau)\le 2u/D$ and $|r(u(\tau))|\le B(u/D)^3$.
Applying Gr\"onwall's inequality to \eqref{eq:delta_eq} gives
\begin{equation}
\label{eq:delta_bound}
|\delta(\ln 2)|\le \exp\Bigl(\int_0^{\ln 2} 2b_0\frac{2u}{D}\,d\tau\Bigr)\,\int_0^{\ln 2} B\Bigl(\frac{u}{D}\Bigr)^3 d\tau
\le \frac{B\ln 2}{D^3}\,u^3.
\end{equation}

\paragraph{Step 3: combine.}
Finally,
$\Sigma(u)-u-2b_0u^2\ln 2 = [w(\ln 2)-u-2b_0u^2\ln 2] + \delta(\ln 2)$.
Combine \eqref{eq:w_trunc_bound} and \eqref{eq:delta_bound} to obtain \eqref{eq:step_expansion_bound} with \eqref{eq:Csigma}.
\end{proof}

\paragraph{Uniform $O(u^3)$ remainder and identification of $b_0$.}
Lemma~\ref{lem:beta_to_step_remainder} reduces the step-scaling remainder to the uniform beta remainder bound \eqref{eq:beta_remainder_assumption}.
This bound is supplied by the analytic dyadic RG theorem:
Proposition~\ref{ass:RG_analytic_bounds} (proved in Theorem~\ref{thm:RG_analytic_bounds}) yields
\[
|\delta_u(u,K)|\le C_u\bigl(|u|^3+|u|\|K\|_j\bigr)
\]
uniformly in the regulator ratio $a/L$ along the dyadic trajectory.
Restricting to the one-dimensional $K=0$ trajectory gives \eqref{eq:beta_remainder_assumption} with $B=C_u$
(after the fixed-scheme identification of $u$ with the gradient-flow coupling).

The universal quadratic coefficient is identified in Proposition~\ref{prop:b0_scheme_invariant}.


\begin{itemize}
\item M.~L\"uscher, ``Properties and uses of the Wilson flow in lattice QCD,'' arXiv:1006.4518 / JHEP 1008 (2010) 071.\newline
Background on the Wilson flow and the tree-level normalization of flowed observables. (Cited for context; the load-bearing coefficient positivity in this submission is obtained internally via Appendix~F.)

\item M.~L\"uscher, P.~Weisz, ``Perturbative analysis of the gradient flow in non-abelian gauge theories,'' arXiv:1101.0963 / JHEP 1102 (2011) 051.\newline
Background on renormalization/UV finiteness of flowed correlators once the underlying 4D theory is renormalized.

\item P.~Fritzsch, A.~Ramos, ``The gradient flow coupling in the Schr\"odinger functional,'' arXiv:1301.4388.
\newline Defines $g^2_{\mathrm{GF}}(L)=\mathcal N^{-1}t^2\langle E(t,x_0)\rangle$ with $t=c^2L^2/8$ and tree-level normalization.
\end{itemize}
